\documentclass{article}
\usepackage{amsmath}
\usepackage{amsfonts}

\begin{document}

\section*{2025 MIT Science Bowl High School Invitational}
\subsection*{Round 14}

\subsection*{TOSS-UP}
\noindent 1) MATH Short Answer Two concentric circles have radii 5 and 10, respectively. Point A is picked on the outer circle, and the two tangent lines to the inner circle passing through A are drawn. In degrees, what is the smaller angle formed by these two lines?
\noindent ANSWER: 60

\subsection*{BONUS}
\noindent 1) MATH Multiple Choice Colin is standing on a 2-dimensional plane at the point (0,1) facing in the positive x direction. He steps forwards 1 unit, and then turns 90 degrees clockwise. He then steps forwards 2 units, turns 90 degrees clockwise, and continues this process, where the distance he steps each turn is equal to the sum of the distances traveled in his two previous steps. Which of the following is the closest to the number of turns Colin has taken by the time he reaches a distance 1000 units from the origin?
\noindent W) 5
\noindent X) 15
\noindent Y) 25
\noindent Z) 35
\noindent ANSWER: X) 15

\subsection*{TOSS-UP}
\noindent 2) EARTH AND SPACE Short Answer Order the following three water masses in order of increasing average density: 1) Mediterranean intermediate water; 2) Antarctic bottom water; 3) North Atlantic deep water.
\noindent ANSWER: 1, 3, 2

\subsection*{BONUS}
\noindent 2) EARTH AND SPACE Short Answer The Pineapple Express is an example of what meteorological system that consist of narrow atmospheric bands rich in water vapor?
\noindent ANSWER: ATMOSPHERIC RIVERS

\subsection*{TOSS-UP}
\noindent 3) PHYSICS Short Answer The ground-state energy of a 1-dimensional ideal quantum harmonic oscillator is 1.5 electron-volts, and the first excited state is 4.5 electron-volts. In electron-volts, what is the energy of the oscillator's second excited state?
\noindent ANSWER: 7.5

\subsection*{BONUS}
\noindent 3) PHYSICS Short Answer Jake has a block on a frictionless incline that is 1 degree above the horizontal, which reaches some terminal velocity due to quadratic air resistance. Serena has an identical block on a frictionless incline that is 2 degrees above the horizontal. When both blocks are at terminal velocity, find the ratio between the power that air resistance dissipates on Serena's block and the power that air resistance dissipates on Jake's block, rounded to two significant figures.
\noindent ANSWER: 2.8

\subsection*{TOSS-UP}
\noindent 4) CHEMISTRY Short Answer Identify all of the following three diatomic molecules that have a net dipole moment: 1) Dihydrogen; 2) Hydrogen deuteride; 3) Dideuterium.
\noindent ANSWER: 2 ONLY

\subsection*{VISUAL BONUS}
\noindent 4) CHEMISTRY Short Answer Your next bonus is visual. A pink-colored complex is formed by dissolving cobalt (II) chloride in water. Answer the following two questions about the electronic structure of the solvated cobalt complex using the image shown: 1) The complex's pink color is caused by a transition between what two molecular orbitals? 2) Using the empty MO diagram provided, what is the average bond order of the cobalt oxygen bond in this complex?
\noindent ANSWER: 1) $t_{2g}$ AND $e_{g}^{z}$ 2) $5/6$

\subsection*{TOSS-UP}
\noindent 5) BIOLOGY Multiple Choice Cobalamin is required for the regeneration of which of the other vitamins?
\noindent W) Niacin
\noindent X) Pantothenic acid
\noindent Y) Pyroxidine
\noindent Z) Folic acid
\noindent ANSWER: Z) FOLIC ACID

\subsection*{VISUAL BONUS}
\noindent 5) BIOLOGY Short Answer Your next bonus is visual. Shown in the left diagram is a 9 mb plasmid with two lox sites facing opposite directions, and some amount of EcoRI sites. You treat the plasmid with EcoRI, cre recombinase, and both cre + EcoRI. After linearizing the products, you run the results on a gel shown on the right. Which of the options in the image best represents the EcoRI cut site configuration on this plasmid?
\noindent ANSWER: C

\subsection*{TOSS-UP}
\noindent 6) ENERGY Multiple Choice The Surendranath Group used a platinum catalyst to facilitate the transfer of hydride anions between molecules. Which of the following molecules would be the best hydride donor with platinum catalysis?
\noindent W) Methane
\noindent X) Ammonia
\noindent Y) Water
\noindent Z) Hydrogen fluoride
\noindent ANSWER: W) METHANE

\subsection*{VISUAL BONUS}
\noindent 6) ENERGY Short Answer Your next bonus is visual. Researchers at the Broad have been studying emerging infectious diseases. Shown in the image is a certain infectious virus predominantly found in West Africa. Answer the following three questions: 1) What is the identity of this virus? 2) What type of fever that involves massive internal bleeding does this virus cause? 3) Would Northern or Southern blotting be more effective for quantifying viral spread?
\noindent ANSWER: 1) EBOLAVIRUS; 2) HEMORRHAGIC FEVER; 3) NORTHERN BLOTTING

\subsection*{TOSS-UP}
\noindent 7) EARTH AND SPACE Multiple Choice While examining a sample of peridotite, Jack is surprised to find a large quartz crystal embedded within it. Which of the following terms best describes the quartz crystal?
\noindent W) Xenoblast
\noindent X) Xenolith
\noindent Y) Phenocryst
\noindent Z) Xenocryst
\noindent ANSWER: Z) XENOCRYST

\subsection*{VISUAL BONUS}
\noindent 7) EARTH AND SPACE Short Answer Your next bonus is visual. Answer the following two questions regarding the Nile River delta: 1) What term best describes the Rossetta and Damietta rivers relative to the Nile River? 2) There are also an abundance of old river channels seen in the delta. What is the name of the phenomenon where a river rapidly abandons a river channel for a new path to sea level?
\noindent ANSWER: 1) DISTRIBUTARIES; 2) RIVER AVULSION (ACCEPT: 2) DELTA SWITCHING)

\subsection*{TOSS-UP}
\noindent 8) ENERGY Multiple Choice The Smith Lab used polymers containing strong carbon nitrogen double bonds to design filtration membranes for hydrocarbons. Which of the following classes of polymers did they use?
\noindent W) Polyamines
\noindent X) Polyamides
\noindent Y) Polyimines
\noindent Z) Polynitriles
\noindent ANSWER: Y) POLYIMINES (READ: POLY - IH - MEENS)

\subsection*{BONUS}
\noindent 8) ENERGY Multiple Choice Physicists from the Frebel Group studied the synthesis of magnesium during star formation. Which of the following stellar events is responsible for producing most of the magnesium in the universe?
\noindent W) Exploding white dwarfs
\noindent X) Exploding massive stars
\noindent Y) Cosmic ray fission
\noindent Z) Merging neutron stars
\noindent ANSWER: X) EXPLODING MASSIVE STARS

\subsection*{TOSS-UP}
\noindent 9) BIOLOGY Short Answer Identify all of the following three synaptic junctions in which acetylcholine is the primary neurotransmitter: 1) Parasympathetic preganglionic; 2) Parasympathetic postganglionic; 3) Sympathetic preganglionic.
\noindent ANSWER: ALL

\subsection*{BONUS}
\noindent 9) BIOLOGY Short Answer Identify all of the following three organisms that underwent a secondary endosymbiosis of a red algae: 1) The causative agent of african sleeping sickness; 2) The organism for which the Anopholes mosquito is the sole vector for; 3) The organisms that cause red tide.
\noindent ANSWER: 2 AND 3

\subsection*{TOSS-UP}
\noindent 10) PHYSICS Multiple Choice A negatively charged particle is initially located a distance d from an infinitely long rod of uniform positive charge density. Which of the following best describes the escape velocity of the particle from the rod?
\noindent W) Proportional to $1/\sqrt{d}$
\noindent X) Proportional to $\ln d$
\noindent Y) Proportional to $\sqrt{\ln(1/d)}$
\noindent Z) Diverges to infinity
\noindent ANSWER: Z) DIVERGES TO INFINITY

\subsection*{BONUS}
\noindent 10) PHYSICS Short Answer A bacterium is swimming in an incompressible fluid. Assuming that the Reynolds number of its motion is much less than one, identify all of the following three statements that are true: 1) Inertial forces dominate over viscous forces; 2) Drag forces are proportional to both speed and viscosity; 3) Gravity and convection are both negligible at the bacterium's scale.
\noindent ANSWER: 2 AND 3

\subsection*{TOSS-UP}
\noindent 11) MATH Multiple Choice The Cauchy distribution is a famous counterexample to the statement that every well-defined distribution must have a well-defined mean and variance. Which of the following best describes the Cauchy distribution?
\noindent W) The distribution of the sum of two normally distributed random variables
\noindent X) The distribution of the difference of two normally distributed random variables
\noindent Y) The distribution of the product of two normally distributed random variables
\noindent Z) The distribution of the quotient of two normally distributed random variables
\noindent ANSWER: Z) THE DISTRIBUTION OF THE QUOTIENT OF TWO NORMALLY DISTRIBUTED RANDOM VARIABLES

\subsection*{VISUAL BONUS}
\noindent 11) MATH Short Answer Your next bonus is visual. A regular tetrahedron ABCD is inscribed in a sphere. Let $A^{\prime}$ be the point diametrically opposite to A on the sphere. In degrees, what is the measure of angle BA'C?
\noindent ANSWER: 90

\subsection*{TOSS-UP}
\noindent 12) CHEMISTRY Multiple Choice Which of the following radical abstraction reactions generally takes place most quickly?
\noindent W) Abstraction of hydrogen from primary carbon by fluorine radical
\noindent X) Abstraction of hydrogen from tertiary carbon by fluorine radical
\noindent Y) Abstraction of hydrogen from primary carbon by chlorine radical
\noindent Z) Abstraction of hydrogen from tertiary carbon by chlorine radical
\noindent ANSWER: X) ABSTRACTION OF HYDROGEN FROM TERTIARY CARBON BY FLUORINE RADICAL

\subsection*{BONUS}
\noindent 12) CHEMISTRY Short Answer Identify all of the following three statements that are true about the potential energy surface of a reaction: 1) Intermediates correspond to local minima on the surface; 2) Transition states correspond to local maxima on the surface; 3) Catalysts alter the surface by lowering the global minima.
\noindent ANSWER: 1 ONLY

\subsection*{TOSS-UP}
\noindent 13) BIOLOGY Short Answer When the ciliary body overproduces vitreous humor, it can lead to what condition in the eye that causes increased intraocular pressure, and can be tested for with a blow test?
\noindent ANSWER: GLAUCOMA

\subsection*{VISUAL BONUS}
\noindent 13) BIOLOGY Short Answer Your next bonus is visual. Shown in the image are three organisms. A is a type of lungfish, and C is a type of salamander. Answer the following two questions: 1) By number, identify all of the following three traits that are shared among these organisms: 1) Amniotic egg; 2) Mineralized skeleton; 3) Jaws; 2) Identify which of these organisms would contain a spiral valve.
\noindent ANSWER: 1) 3 ONLY; 2) B

\subsection*{TOSS-UP}
\noindent 14) CHEMISTRY Multiple Choice Which of the following substituted benzoic acids has the highest acid dissociation constant in aqueous solution?
\noindent W) Ortho nitro benzoic acid
\noindent X) Ortho methoxy benzoic acid
\noindent Y) Meta nitro benzoic acid
\noindent Z) Meta methoxy benzoic acid
\noindent ANSWER: W) ORTHO NITRO BENZOIC ACID

\subsection*{BONUS}
\noindent 14) CHEMISTRY Short Answer Rank the following four molecular orbitals by increasing energy: 1) HOMO of ethylene; 2) LUMO of ethylene; 3) HOMO of 1,3-butadiene; 4) LUMO of 1,3-butadiene.
\noindent ANSWER: 1, 3, 4, 2

\subsection*{TOSS-UP}
\noindent 15) MATH Multiple Choice Which of the following is not a property of an invertible $n\times n$ matrix A?
\noindent W) The determinant of A is nonzero
\noindent X) The rows of A are linearly independent
\noindent Y) The transpose of A is an invertible matrix.
\noindent Z) The kernel of A has dimension n
\noindent ANSWER: Z) THE KERNEL OF A HAS DIMENSION n

\subsection*{BONUS}
\noindent 15) MATH Multiple Choice The polynomial $x^{4}-2x^{2}=a$, where a is a real number, has n distinct real roots. What is the sum of all possible values of n?
\noindent W) 2
\noindent X) 4
\noindent Y) 6
\noindent Z) 9
\noindent ANSWER: Z) 9

\subsection*{TOSS-UP}
\noindent 16) EARTH AND SPACE Multiple Choice Which of the following characteristics of a main-sequence star is negatively correlated with its surface temperature?
\noindent W) Radius
\noindent X) Luminosity
\noindent Y) Escape velocity
\noindent Z) Opacity
\noindent ANSWER: Z) OPACITY

\subsection*{BONUS}
\noindent 16) EARTH AND SPACE Short Answer Identify all of the following three observations that the current Lambda-CDM model of the universe struggles to explain: 1) Galaxies in the early universe are too massive; 2) CDM simulations predict larger numbers of small dark matter halos than are actually observed; 3) Galaxies outside the Local Group are moving away faster than predicted.
\noindent ANSWER: ALL

\subsection*{TOSS-UP}
\noindent 17) ENERGY Multiple Choice The Kaiser Group recently discovered that microscopic non-rotating black holes could carry a significant color charge. These black holes are analagous to which of the following solutions to the Einstein field equations?
\noindent W) Kerr
\noindent X) Kerr-Newman
\noindent Y) Schwarzschild
\noindent Z) Reissner-Nordström
\noindent ANSWER: Z) REISSNER-NORDSTRÖM

\subsection*{BONUS}
\noindent 17) ENERGY Short Answer Physicists at MIT supported by the DOE recently reported a study on the final particle distribution resulting from collisions between protons. Identify all of the following three particles that could be directly produced by the collision between two protons: 1) Pion; 2) Neutrino; 3) Kaon.
\noindent ANSWER: 1 AND 3

\subsection*{TOSS-UP}
\noindent 18) PHYSICS Short Answer Lydia is traveling at a constant speed of 0.6c away from Nathan. Rounded to one significant figure and expressed in scientific notation in meters squared per seconds squared, what is the magnitude of Lydia's four-velocity in Nathan's reference frame?
\noindent ANSWER: $9\times10^{16}$ (ACCEPT: -9x $10^{16}$)

\subsection*{BONUS}
\noindent 18) PHYSICS Multiple Choice An isolated quantum particle is in a superposition of the ground-state and first excited state in a one-dimensional potential well. Which of the following describes the probability distribution of the particle's energy after a long time?
\noindent W) Equal to the ground-state energy with high probability
\noindent X) Equal to the first excited state energy with high probability
\noindent Y) The probability distribution oscillates between the ground-state and first excited state energies
\noindent Z) The probability distribution remains the same
\noindent ANSWER: Z) THE PROBABILITY DISTRIBUTION REMAINS THE SAME

\subsection*{TOSS-UP}
\noindent 19) MATH Short Answer The sum of the even divisors of a positive integer n, including n itself, is 30 times the sum of the odd divisors of n. What is the largest power of 2 that is a divisor of n?
\noindent ANSWER: 16

\subsection*{BONUS}
\noindent 19) MATH Short Answer Katie has a finite set of real numbers with mean 1. Identify all of the following three quantities that CANNOT decrease when every number in the set is squared: 1) Mean; 2) Standard deviation; 3) Range.
\noindent ANSWER: 1 ONLY

\subsection*{TOSS-UP}
\noindent 20) CHEMISTRY Short Answer To the nearest inverse meter, what is the wavenumber of light emitted by the spin-flip transition of a hydrogen atom's lone electron?
\noindent ANSWER: 5

\subsection*{BONUS}
\noindent 20) CHEMISTRY Short Answer Identify all of the following three substituents on a carbocation that would destabilize the adjacent positive formal charge: 1) Phenyl; 2) Chloro; 3) Hydroxyl.
\noindent ANSWER: 2 ONLY

\subsection*{TOSS-UP}
\noindent 21) PHYSICS Multiple Choice According to the theory of general relativity, which of the following observers has nonzero proper acceleration?
\noindent W) A student seated in a classroom
\noindent X) An astronaut orbiting the earth
\noindent Y) A skydiver free falling from a plane with no air resistance
\noindent Z) There is no universal notion of acceleration
\noindent ANSWER: W) A STUDENT SEATED IN A CLASSROOM

\subsection*{VISUAL BONUS}
\noindent 21) PHYSICS Short Answer Your next bonus is visual. Cyan is considering Circuit A, shown in the image, which consists of an independent voltage source, some resistors, and an inductor. From Thevenin's theorem, she knows that this circuit is equivalent to Circuit B. Answer the following two questions: 1) In volts, what is the Thevenin voltage V? 2) In kiloohms, what is the Thevenin resistance R?
\noindent ANSWER: 1) 4 VOLTS; 2) 4 KILOOHMS

\subsection*{TOSS-UP}
\noindent 22) EARTH AND SPACE Short Answer Some prograde rotating asteroids in the solar system are observed to have their semimajor axis slowly increase over time, thus spiraling away from the sun. This is due to the illumination of the sun causing anisotropic emission of photons, which is known by what effect?
\noindent ANSWER: YARKOVSKY EFFECT

\subsection*{VISUAL BONUS}
\noindent 22) EARTH AND SPACE Short Answer Your next bonus is visual. Shown in the image is a graph of radial velocity over time for each star in a binary system with an unknown inclination relative to Earth. It is known that Star A has mass 5 solar masses and constant distance 8 astronomical units from the system's barycenter. What is the period of this system, in years?
\noindent ANSWER: 56

\subsection*{TOSS-UP}
\noindent 23) BIOLOGY Short Answer Smoltification in salmon, which is the physiological adaptation to living in salty water, is induced by the release of what hormone, which is also produced by the zona fasciculata in humans?
\noindent ANSWER: CORTISOL

\subsection*{BONUS}
\noindent 23) BIOLOGY Short Answer You are investigating different structures cells form during cytokinesis. You have stained cells undergoing cytokinesis with DAPI and an antibody against tubulin. Which of the following groups of organisms would be distinguishable from charophytes? 1) Chlorophytes; 2) Land plants; 3) Red algae.
\noindent ANSWER: 1 AND 3

\subsection*{TOSS-UP}
\noindent 24) ENERGY Short Answer The Fu Group used fermionic neural networks to minimize the ground-state energy of an exotic state of matter. What method from quantum mechanics do fermionic neural networks use to approximate a computable form of the system's ground-state energy?
\noindent ANSWER: VARIATIONAL (ACCEPT: VARIATIONAL PRINCIPLE)

\subsection*{BONUS}
\noindent 24) ENERGY Short Answer Researchers at the Coley Group used machine learning to identify top-performing compounds from chemical libraries. Identify all of the following three assumptions that must have been made for the researchers to use an ordinary least squares model: 1) Regressors and errors are uncorrelated; 2) Regressors are independent; 3) The variance of each error term is identical.
\noindent ANSWER: ALL

\subsection*{TOSS-UP}
\noindent 25) CHEMISTRY Short Answer If the spin-only magnetic moment of a single-center transition metal complex is approximately 3.9 bohr magnetons, how many unpaired electrons does the complex's metal center have?
\noindent ANSWER: 3

\subsection*{BONUS}
\noindent 25) CHEMISTRY Short Answer A 1-liter sample of an ideal gas expands isothermally at 150 kelvin to a final volume of 4 liters, and the entropy change of the gas is measured to be S. If the same 1-liter sample of gas instead expands isothermally at 300 kelvin to a final volume of 8 liters, then in terms of S, what is the new entropy change of the gas?
\noindent ANSWER: 1.5S

\end{document}